%%==============================================================================
%% 中文摘要 / 英文摘要
%%==============================================================================

\begin{cabstract}
随着智能制造的快速发展，人机协作已成为现代工业生产的重要模式。然而，在动态、
非结构化的协作环境中，如何对人与机器人的任务进行实时、高效的规划，仍然是一个
具有挑战性的问题。本文以工业人机协作场景为研究对象，从任务建模、动态规划与
策略优化三个方面开展研究，主要工作如下：

（1）针对人机协作任务的不确定性，建立了基于马尔可夫决策过程的任务规划模型，
将人的行为意图作为环境状态的一部分进行建模，提高了任务分配的合理性。

（2）针对动态环境下的实时规划需求，提出了一种基于深度强化学习的动态任务规划
方法。通过设计兼顾效率与安全的奖励函数，使智能体能够在协作过程中自适应地调整
任务分配策略。

（3）搭建了人机协作实验平台，在多种工况下验证了所提方法的有效性。实验结果
表明，与传统方法相比，本文方法在任务完成时间和协作安全性方面均具有明显优势。

\ckeywords{人机协作；深度强化学习；动态任务规划；智能制造；马尔可夫决策过程}
\end{cabstract}

\begin{eabstract}
With the rapid development of intelligent manufacturing, human-robot collaboration
(HRC) has become an important paradigm in modern industrial production. However,
how to plan the tasks of humans and robots in real time and efficiently in a dynamic
and unstructured collaborative environment remains a challenging problem. This thesis
takes the industrial HRC scenario as the research object and carries out research from
three aspects: task modeling, dynamic planning, and policy optimization. The main
work is as follows:

(1) To address the uncertainty of HRC tasks, a task planning model based on the Markov
decision process is established, in which human behavioral intention is modeled as part
of the environment state, improving the rationality of task allocation.

(2) To meet the real-time planning requirements in dynamic environments, a dynamic
task planning method based on deep reinforcement learning is proposed. By designing a
reward function that balances efficiency and safety, the agent can adaptively adjust the
task allocation policy during collaboration.

(3) An HRC experimental platform is built, and the effectiveness of the proposed method
is verified under various working conditions. Experimental results show that, compared
with traditional methods, the proposed method has obvious advantages in task completion
time and collaboration safety.

\ekeywords{Human-Robot Collaboration; Deep Reinforcement Learning; Dynamic Task
Planning; Intelligent Manufacturing; Markov Decision Process}
\end{eabstract}
